\documentclass[11pt,a4paper]{article}
% Shared preamble for RuPhO English problem statements (match T1 2026 style).
% Use: \input{latex/rupho_en_preamble.tex} after \documentclass.
\usepackage[margin=1in]{geometry}
\usepackage{amsmath,amssymb}
\usepackage{graphicx}
\usepackage{float}
\usepackage{enumitem}
\usepackage{microtype}
\usepackage{caption}
\usepackage{tabularx}
\usepackage{hyperref}

\captionsetup{font=small,labelfont=bf,skip=6pt}

\setlist[itemize]{leftmargin=1.2cm}
\setlist[enumerate]{leftmargin=1.2cm}

% One outer box per subproblem: [id | statement | points] (no inner rules)
\newcommand{\problembox}[3]{%
  \par\addvspace{0.42em}%
  \noindent
  \setlength{\fboxsep}{7.5pt}%
  \setlength{\fboxrule}{0.95pt}%
  \fbox{%
    \begin{minipage}{\dimexpr\linewidth-2\fboxsep-2\fboxrule\relax}
    \begin{tabularx}{\linewidth}{@{}>{\raggedright\arraybackslash\bfseries}p{2.1em} @{\hspace{0.9em}} >{\raggedright\arraybackslash}X @{\hspace{0.9em}} >{\raggedleft\arraybackslash\bfseries}p{2.4em}@{}}
    #1 & #2 & #3
    \end{tabularx}
    \end{minipage}%
  }%
  \par\addvspace{0.32em}%
}
\title{\textbf{Electrode-electrolyte interface --- Solution}}
\author{\normalsize X23 (2023) --- English translation}
\date{}
\begin{document}
\maketitle
\problembox{A1}{Measure the inner diameter of the syringe $D_0$ and the distance $L_0$ between marks $0~\text{}$ and $20~\text{}$.\\ Answer: $L_0 = 65 \: mm$}{0.02}
\problembox{A2}{Make resistance measurements $R$ corresponding to the positions of the syringe $V=2,4,\dots,20~\text{}$.}{0.40}
\problembox{A3}{Plot $R$ versus $V$. Draw a smoothing straight line $R=R_1 + \alpha_1 V$. Calculate the values of parameters $R_1$ and $\alpha_1$.\\ Answer: $R_1 = 4.3 \: \Omega$}{0.40}
\begin{figure}[H]
    \centering
    \includegraphics[width=0.72\textwidth]{E1_S_EN_assets/fig1.jpg}
\end{figure}
\problembox{B1}{For each of the prepared solutions, carry out a series of $4$-ex measurements $R$ from $V$.}{0.80}
\begin{figure}[H]
    \centering
    \includegraphics[width=0.72\textwidth]{E1_S_EN_assets/fig2.jpg}
\end{figure}
\problembox{B2}{Using least squares method for each solution, calculate the parameters $R_1$ and $\alpha_1$ by analogy with point $\textbf{3}$. Present the results in the form of a table.\\ c 0.07 0.1 0.13 0.16 0.19 $\alpha_1,~\frac{\text{\Omega}}{cm^3}$ 1.15 0.80 0.66 0.58 0.56 $R_1,~\text{\Omega}$ 4.3 5.0 4.9 4.5 3.7}{0.12}
\problembox{B3}{An external constant field $E$ is created in the medium. What is the steady-state velocity of the particle $v_{}$? Write the answer using $E,k,q$.\\ Let's write down Newton's II law: \[ Eq = kv_\text{y} \]\\ Answer: \[ v_\text{y} = \frac{Eq}{k} \]}{0.06}
\problembox{B4}{What is the steady-state current $I$ if a potential difference $U$ is applied to the electrodes? Neglect any surface effects (eg contact potential difference and contact resistance). Write your answer using $n,S,L,e,k_{\rm Na^+}, k_{\rm Cl^-}$ and $U$.\\ By definition $I = \frac{dq}{dt}$, in this case, a flow of particles with an area $S$ and a charge $q$ moving at a constant speed $v$ forms the following current: \[ I = \frac{dq}{dt}= nvSq  \tag{1}\] Then the current in our case consists of two ion flows. The ion concentrations are equal from the condition of electrical neutrality: \[ I = nSe \frac{Ee}{k_{\rm Na^+}} + nSe \frac{Ee}{k_{\rm Cl^-}}.\] The field $E$ is expressed through the potential difference as $U/L$.\\ Answer: \[ I = \frac{nSe^2U}{L} \left( \frac{1}{k_{\rm Na^+}}+\frac{1}{k_{\rm Cl^-}} \right)\]}{0.15}
\problembox{B5}{Let a potential difference $U(t) = U \cos \omega t$ be applied to the bases, how will the current $I(t)$ depend on time? Again, ignore any surface effects. Write your answer using $n,S,L,e,\omega,m_{\rm Na^+}, m_{\rm Cl^-},k_{\rm Na^+}, k_{\rm Cl^-}, U$ and $t$.\\ We will solve the problem using the method of complex amplitudes: \[ U \cos \omega t = \mathfrak{Re} \, U e^{i \omega t}. \] Then the complex electric field $E$ inside the medium is: \[ \tilde{E} = \frac{U}{L} e^{i \omega t}. \] We write Newton’s II law for a particle in an alternating field $Ee^{i\omega t}$: \[ m\dot{v} + kv = Eq e^{i \omega t}.\] Our own solutions are exponentially decaying, so we are only interested in a particular solution of the form $v=\tilde{B} e^{i \omega t}$: \[ i\omega m \tilde{B} + k\tilde{B} = Eq \quad \Rightarrow \quad \tilde{B} = \frac{Eq}{i \omega m + k} = \frac{Eq}{k} \frac{1}{1 + i\omega\frac{m}{k}} = \frac{Eq}{k} \frac{1}{\sqrt{1+\frac{\omega^2 m^2}{k^2} }}e^{-i\arctan \frac{\omega m}{k}}.\] Complex Let's recalculate the current $\tilde{I}$ taking into account equation (1): \[ \tilde{I} = \frac{nSe^2 U}{Lk_{\rm Na^+}} \frac{1}{\sqrt{1+\frac{\omega^2 m_{\rm Na^+}^2}{k_{\rm Na^+}^2} }} e^{-i\arctan \frac{\omega m_{\rm Na^+}}{k_{\rm Na^+}} +i \omega t} + \frac{nSe^2 U}{Lk_{\rm Cl^-}} \frac{1}{\sqrt{1+\frac{\omega^2 m_{\rm Cl^-}^2}{k_{\rm Cl^-}^2} }} e^{-i\arctan \frac{\omega m_{\rm Cl^-}}{k_{\rm Cl^-}}+i \omega t}\] Let's move from the complex current by taking the real part.\\ Answer: \[ I = \frac{nSe^2 U}{Lk_{\rm Na^+}} \frac{1}{\sqrt{1+\frac{\omega^2 m_{\rm Na^+}^2}{k_{\rm Na^+}^2} }} \cos \left( -\arctan \frac{\omega m_{\rm Na^+}}{k_{\rm Na^+}} +\omega t \right) + \frac{nSe^2 U}{Lk_{\rm Cl^-}} \frac{1}{\sqrt{1+\frac{\omega^2 m_{\rm Cl^-}^2}{k_{\rm Cl^-}^2} }} \cos \left( -\arctan \frac{\omega m_{\rm Cl^-}}{k_{\rm Cl^-}} +\omega t \right) \]}{0.65}
\problembox{B6}{Now let's take into account surface effects in the form of a resistor $R_\text{p}$, which is connected in series to the solution. Write down the formula for the impedance $\tilde{Z}(\omega)$ of a system consisting of a solution $\rm NaCl$ and a vessel. Write your answer using $n,S,L,e,m_{\rm Na^+}, m_{\rm Cl^-},\omega,k_{\rm Na^+}, k_{\rm Cl^-}$ and $R_\text{p}$.\\ The impedance of one type of ion is obviously expressed from the already obtained equations \[ \tilde{Z} = \frac{Lk}{nSq^2}  \left( 1+  i \omega \frac{m}{k} \right),\], while the currents add up at the same voltages, which means that impedances $\rm Na^+$ and $\rm Cl^-$ are “connected” in parallel, and there is also a resistor $R_\text{p}$ in series with them.\\ Answer: \[ \tilde{Z}(\omega) = R_\text{p} + \frac{\frac{k_{\rm Na^+}L}{nSe^2} \left(1 +  \frac{i \omega m_{\rm Na^+}}{k_{\rm Na^+}} \right) \cdot \frac{k_{\rm Cl^-}L}{nSe^2} \left(1 +  \frac{i \omega m_{\rm Cl^-}}{k_{\rm Cl^-}} \right)} {\frac{k_{\rm Na^+}L}{nSe^2} \left(1 +  \frac{i \omega m_{\rm Na^+}}{k_{\rm Na^+}} \right) + \frac{k_{\rm Cl^-}L}{nSe^2} \left(1 +  \frac{i \omega m_{\rm Cl^-}}{k_{\rm Cl^-}} \right)}\]}{0.25}
\problembox{B7}{What is the theoretical system impedance $\tilde{Z}$ if the phase difference between the voltages is negligible? Express your answer in terms of $n,S,L,e,k_{\rm Na^+}, k_{\rm Cl^-}$ and $R_\text{p}$.\\ If the phase difference is negligible, then the complex parts can be neglected and the behavior of the system corresponds to the behavior at a constant voltage.\\ Answer: \[ \tilde{Z} = R_\text{p} + \frac{L}{nSe^2} \frac{k_{\rm Na^+} \cdot k_{\rm Cl^-}}{k_{\rm Na^+} + k_{\rm Cl^-}} \]}{0.10}
\problembox{B8}{Using the linearized graph of $\alpha_1$ against $c$ obtain the value $\alpha$.\\ Answer: \[ \alpha_1 = \frac{dZ}{dL} \cdot \frac{1}{S} \propto \frac{k}{n} \propto n^{\alpha -1} \] From graph $\alpha_1 = 0.3$.}{0.40}
\begin{figure}[H]
    \centering
    \includegraphics[width=0.72\textwidth]{E1_S_EN_assets/fig3.jpg}
\end{figure}
\problembox{B9}{Obtain the theoretical dependence $|\tilde{Z}|^2(\omega)$ in the approximation $\omega m_{\rm Na^+} \gg k_{\rm Na^+}$, $\omega m_{\rm Cl^-} \gg k_{\rm Cl^-}$. Write your answer using $n,S,L,e,m_{\rm Na^+}, m_{\rm Cl^-},\omega,k_{\rm Na^+}, k_{\rm Cl^-}$ and $R_\text{p}$.\\ Let's simplify the formula for impedance from B6: \[ \tilde{Z} = R_\text{p} + \frac{L}{nSe^2} \frac{(k_{\rm Na^+} + i \omega m_{\rm Na^+})(k_{\rm Cl^-} + i \omega m_{\rm Cl^-})}{(k_{\rm Na^+} + k_{\rm Cl^-}) + i \omega (m_{\rm Na^+}+m_{\rm Cl^-})} \] The first approximation occurs at the multiplication stage: \[ \tilde{Z} = R_\text{p} + \frac{L}{nSe^2} \frac{ -\omega^2 m_{\rm Na^+} m_{\rm Cl^-}+ i \omega( m_{\rm Na^+}k_{\rm Cl^-} +m_{\rm Cl^-}k_{\rm Na^+}) }{(k_{\rm Na^+} + k_{\rm Cl^-}) + i \omega (m_{\rm Na^+}+m_{\rm Cl^-})} \] Multiply the numerator and denominator by the complex conjugate of the denominator and neglect again: \[ \tilde{Z} = R_\text{p} + \frac{L}{nSe^2} \frac{ - \omega^2 m_{\rm Na^+} m_{\rm Cl^-} (k_{\rm Na^+} + k_{\rm Cl^-})+ i\omega^3 m_{\rm Na^+}  m_{\rm Cl^-} (  m_{\rm Na^+} + m_{\rm Cl^-})   + \omega^2 ( m_{\rm Na^+}k_{\rm Cl^-} +m_{\rm Cl^-}k_{\rm Na^+}) (m_{\rm Na^+} + m_{\rm Cl^-}) }{ \omega ^2(m_{\rm Na^+}+m_{\rm Cl^-})^2} \]\\ Answer: \[ |\tilde{Z}|^2 = \left( R_\text{p} + \frac{L}{nSe^2} \frac{m_{\rm Na^+}^2 k_{\rm Cl^-} + m_{\rm Cl^+}^2 k_{\rm Na^+} }{(m_{\rm Na^+} + m_{\rm Cl^-})^2} \right)^2 + \omega^2 \left( \frac{L}{nSe^2} \frac{m_{\rm Na^+} m_{\rm Cl^-}}{m_{\rm Na^+} + m_{\rm Cl^-}} \right)^2 \]}{0.40}
\problembox{B10}{For a solution with salt concentration $c\approx 0.20$, measure the impedance modulus $|Z|$ versus frequency $f$. The series must consist of at least 7 different $f$ in the range $500~\text{kHz} \leq f \leq 10~\text{MHz}$.}{0.35}
\problembox{B11}{Plot a linearized graph of $|Z|$ versus $f$. From the graph, determine the experimental value $\mu_\text{eff}=\frac{\mu_{\rm Na} \cdot \mu_{\rm Cl}}{\mu_{\rm Na} + \mu_{\rm Cl}}$.\\ According to the last theoretical point, the graph of $|\tilde{Z}|^2$ versus $f^2$ is linear.\\ Answer: From the graph, the slope coefficient is $61\cdot 10^{-12}~\frac{\text{\Omega}^2}{\text{Hz}^2}$, and from theory it is $\left( \frac{2\pi L m_\text{eff}}{nSe^2}\right)^2$. $\mu_{eff} = 170 \cdot 10^{5} \: \frac{kg}{mol}$}{0.50}
\begin{figure}[H]
    \centering
    \includegraphics[width=0.72\textwidth]{E1_S_EN_assets/fig4.jpg}
\end{figure}
\problembox{B12}{Determine $N$ based on your experimental data.\\ Let's write the formula based on the model described in the condition: \[ \mu_\text{eff} = \frac{(\mu_{\rm Na}+ N \mu_{\rm H_2 O})\cdot (\mu_{\rm Cl}+ N \mu_{\rm H_2 O})}{\mu_{\rm Na}+ N \mu_{\rm H_2 O} + \mu_{\rm Cl}+ N \mu_{\rm H_2 O}}\] Let's transform it to a quadratic equation for $N$: \[ N^2 + N \left( \frac{\mu_{\rm Na}+\mu_{\rm Cl}- 2 \mu_\text{eff} }{\mu_{\rm H_2O}} \right) + \frac{\mu_{\rm Na}\mu_{\rm Cl} - \mu_\text{eff} (\mu_{\rm Na} + \mu_{\rm Cl})}{\mu_{\rm H_2O}^2}=0\] Due to the fact that $\mu_\text{eff}$ is very larger, in fact, $N=\frac{2\mu_\text{eff}}{\mu_{\rm H_2O}}=1.8\cdot 10^{7}$. It can be concluded that the measured inductance is in no way related to the theoretically considered effect. Meanwhile, the inductance of a turn with a current radius of $10~\text{cm}$ can be estimated as $10^{-6}~\text{H}$, which is of the same order as the root of the slope coefficient.\\ Answer: \[ N=1.8\cdot 10^{7} \]}{0.15}
\problembox{C1}{In the frequency range $5~\text{Hz} \leq f \leq 1~\text{kHz}$, measure the syringe impedance modulus $|\tilde{Z}|$ and the time $\Delta t$ by which the current peak precedes the voltage peak. Take at least 15 measurements.}{1.50}
\problembox{C2}{Based on your data, plot a Nyquist plot of $\mathfrak{Im} \tilde{Z}_\text{p}$ from $\mathfrak{Re} \tilde{Z}_\text{p}$ for surface effects. Indicate on the diagram in the form of an arrow the direction of transition from point to point, corresponding to the increase $\omega$.}{0.60}
\begin{figure}[H]
    \centering
    \includegraphics[width=0.72\textwidth]{E1_S_EN_assets/fig5.jpg}
\end{figure}
\problembox{C3}{Qualitatively draw the Nyquist diagram for a parallel-plate capacitor.}{0.10}
\begin{figure}[H]
    \centering
    \includegraphics[width=0.72\textwidth]{E1_S_EN_assets/fig6.jpg}
\end{figure}
\problembox{C4}{Write down the expressions for $C_l$ and $R_l$ in terms of $r,\sigma,\rho$ and $\xi$.\\ Answer: Capacitance $C_l$ is the capacitance per unit length of the surface of the cylinder, so \[C_l=2\pi r \sigma.\] Resistance $R_l$ is the resistance per unit length of water around one cylinder, so $R_l = \rho / S_1$, where $S_1$ is the area of water per cylinder, so \[ R_l = \frac{\rho}{\frac{1}{\xi}-\pi r^2}. \]}{0.20}
\problembox{C5}{Write the complex current $\tilde{I}(x)$ in terms of the derivative of the complex potential $\frac{d\tilde{\varphi}(x)}{dx}$ and $R_l$.\\ Let's write down the voltage drop per unit length of the conductor: \[ -R_ldx \, \tilde{I}  = \tilde{\varphi}(x+dx) - \tilde{\varphi}(x)  \]\\ Answer: \[ \tilde{I}(x) = -\frac{1}{R_l}\frac{d\tilde{\varphi}(x)}{dx} \]}{0.10}
\problembox{C6}{Relate the derivative of the complex current $\frac{d\tilde{I}(x)}{dx}$ and the complex potential $\tilde{\varphi}(x)$ via $C_l$.\\ Current through unit length of capacitor \[ \frac{\tilde{I}(x+dx) - \tilde{I}(x)}{i\omega C_l dx} = -\tilde{\varphi} (x) \]\\ Answer: \[ \frac{d \tilde{I}(x)}{d x} = -i \omega C_l \tilde{\varphi}(x) \]}{0.10}
\problembox{C7}{Show that $\tilde{I}(0) = 0$.\\ If the current is not zero, then it flows through an infinitesimal capacitance, which means $\varphi(0) \to \infty$.}{0.20}
\problembox{C8}{Solve the system of differential equations obtained in steps $\textbf{C5}$ and $\textbf{C6}$ and obtain the impedance $\tilde{Z}_\text{}(\omega)$ of one cylinder and the electrolyte surrounding it. Neglect the capacity of the base of the cylinder and the capacity of the flat surface on which the cylinders stand. Express your answer using $C_l,R_l,l$ and $\omega$.\\ We obtain a differential equation on $\tilde{\varphi}(x)$: \[ \tilde{\varphi}'' - i \omega C_l R_l \tilde{\varphi} = 0 \] This is a homogeneous linear differential equation of the second order. Let's find our own numbers: \[ \lambda^2 = i \omega C_l R_l. \] \[ \lambda_1 = e^{\frac{i \frac{\pi}{2}}{2}} \sqrt{\omega C_l R_l} = (1+i) \sqrt{\frac{\omega R_l C_l}{2}} \] \[ \lambda_2 = e^{\frac{i \frac{\pi}{2}+2\pi i}{2}} \sqrt{\omega C_l R_l} = -(1+i) \sqrt{\frac{\omega R_l C_l}{2}} \] Let's denote $\sqrt{\frac{\omega C_l R_l}{2}}$ as $\zeta$. Then \[ \tilde{\varphi}(x) = A e^{i \zeta x} e^{\zeta x} + B e^{-i \zeta x} e^{-\zeta x}.\] Boundary conditions: $\tilde{\varphi}(l)=Ue^{i \omega t}$, $\tilde{\varphi}'(0) = 0$. \[ \tilde{\varphi}'(0) = A(1+i) \zeta - B(1+i)\zeta = 0 \quad \Rightarrow \quad B = A,\] \[ \tilde{\varphi}(l) = A \left( e^{i \zeta l} e^{\zeta l} + e^{-i \zeta l} e^{-\zeta l} \right) = 2A \cosh(\zeta l + i \zeta l) = U e^{i \omega t} \] As a result, the solution for the potential looks like this: \[ \tilde{\varphi}(x) = U e^{i \omega t} \frac{ \cosh(\zeta x+i \zeta x)}{\cosh(\zeta l+ i \zeta l)}\] Let's find the current $\tilde{I}(l)$: \[ \tilde{I}(l) = -\frac{1}{R_l} \tilde{\varphi}'(l) = -U e^{i \omega t} \frac{\zeta (1+i)}{R_l} \tanh(\zeta l+ i \zeta l), \] then the impedance (the minus arises due to the coordination of the positive direction of the current and the decrease in impedance): \[\tilde{Z}_\text{} = \frac{\tilde{\varphi}(l)}{-\tilde{I}(l)}= \frac{R_l}{\zeta (1+i) \tanh (\zeta l + i \zeta l) } = \sqrt{\frac{R_l}{2\omega C_l}} \frac{1-i}{\tanh \left( (1+i) \sqrt{\frac{\omega C_l R_l}{2}} l \right)} \] Let's expand $\tanh$: \[\tilde{Z}_\text{} = \sqrt{\frac{R_l}{\omega C_l}} \frac{1 - i\tanh \sqrt{\frac{\omega C_l R_l}{2}} l \cdot \tan \sqrt{\frac{\omega C_l R_l}{2}} l}{\tanh \sqrt{\frac{\omega C_l R_l}{2}} l - i\tan \sqrt{\frac{\omega C_l R_l}{2}} l}  e^{-i\pi/4}. \]\\ Answer: \[\tilde{Z}_\text{} = \sqrt{\frac{R_l}{\omega C_l}} \frac{1 + i\tanh \sqrt{\frac{\omega C_l R_l}{2}} l \cdot \tan \sqrt{\frac{\omega C_l R_l}{2}} l} {\tanh \sqrt{\frac{\omega C_l R_l}{2}} l + i\tan \sqrt{\frac{\omega C_l R_l}{2}} l}  e^{-i\pi/4}. \]}{1.25}
\problembox{C9}{Obtain the impedance $\tilde{Z}_\text{}(\omega)$ to the approximation $\sqrt{\omega C_l R_l l^2} \gg 1$. Express your answer using $C_l,R_l,l$ and $\omega$. The following theoretical points are also solved in this approximation.\\ Let's multiply the numerator and denominator by $\cosh \zeta l \cdot \cos \zeta l$: \[\tilde{Z}_\text{} = \sqrt{\frac{R_l}{\omega C_l}} \frac{\cosh \zeta l \cos \zeta l + i \sin \zeta l \sinh \zeta l} {\cos \zeta l \sinh \zeta l + i \cosh \zeta l \sin \zeta l} e^{-i\pi/4}, \] at $\zeta l \gg 1$ the hyperbolic functions $\sinh \zeta l$ and $\cosh \zeta l$ are equal to the exponent $\frac{e^{\zeta l}}{2}$, therefore \[ \tilde{Z}_\text{} \simeq \sqrt{\frac{R_l}{\omega C_l}} \frac{\cos \zeta l e^{\zeta l} + i \sin \zeta l e^{\zeta l}}{\cos \zeta l e^{\zeta l} + i \sin \zeta l e^{\zeta l}}e^{-i\pi/4} =  \sqrt{\frac{R_l}{\omega C_l}}e^{-i\pi/4} \]\\ Answer: \[ \tilde{Z}_\text{} =  \sqrt{\frac{R_l}{\omega C_l}}e^{-i\pi/4} \]}{0.20}
\problembox{C10}{What is the impedance $\tilde{Z}_\text{p}(\omega)$ of one surface of the cylinder in contact with the electrolyte? Surface area $S$. Express your answer using $r,\sigma,\rho,\xi,S$ and $\omega$. Draw a qualitative Nyquist diagram $\tilde{Z}_\text{p}(\omega)$. The theoretical and experimental Nyquist diagrams should be identical in shape, but some of their characteristics may differ slightly.\\ Let's combine the results of the previous paragraphs. \[ \tilde{Z}_\text{p} = \frac{1}{\xi S} \sqrt{ \frac{\rho}{2\pi r \sigma \omega \left(\frac{1}{\xi} - \pi r^2 \right)} } e^{-i\pi/4}\]}{0.20}
\begin{figure}[H]
    \centering
    \includegraphics[width=0.72\textwidth]{E1_S_EN_assets/fig7.jpg}
\end{figure}
\problembox{11}{Plot a linearized graph of $|\tilde{Z}_\text{p}|$ versus $\omega$.}{0.70}
\begin{figure}[H]
    \centering
    \includegraphics[width=0.72\textwidth]{E1_S_EN_assets/fig8.jpg}
\end{figure}
\problembox{C12}{Define $\sigma$.\\ Answer: $\sigma_0 = 6.9 \: \frac{\mu F}{cm^2}$}{0.10}
\vspace{1em}
\noindent\textit{Source:} \url{https://pho.rs/p/3156/s}
\end{document}